\emph{Attribution rule and nulls.} The 5-minute atom is \emph{our} proxy for one attack action, so
\cref{apptab:semblur} counts instead how many distinct red-team steps each issued alert merges, using
LSPR23's own task record (288 tasks, 295 timestamped submissions,
2023-03-09 07:03--15:08 UTC), an attacker-defined decomposition fixed before the dataset was
published. A step is attributed to an alert when its submission falls in the alert's bucket span
\emph{and} its task's declared network segments meet the alert's, or one of its compromise-report IPs
is an alert endpoint. The link is the segment vocabulary --- the task record's \texttt{bt\_baf\_int}
is the dataset's own per-flow \texttt{baf\_int} --- corroborated against the machine-readable
compromise reports (36/39 IPv4 appear as a flow endpoint;
83 agree, 0 disagree). Two nulls are run. The \textbf{label} null permutes which
task each submission belongs to, holding the times fixed: it asks whether \emph{which} task acted is
informative, and goes degenerate at the daily bucket where every submission shares one bucket. The
\textbf{shift} null rotates the whole red-team timeline against the stream, preserving task identity,
ordering, attributes and inter-arrival structure: it asks the sharper question, whether the alignment
\emph{in time} is informative, and does not degenerate.

\emph{Reading.} The result is agreement in direction without discrimination. The external count and
the proxy rise together (panel (b) of \cref{apptab:semblur}), so the resolution cost survives
substituting an attacker-defined unit for ours; but the count clears the label null at only $6$ of
$19$ cells and the shift null at only $1$, so the rise is driven by a wider bucket capturing more
submissions rather than by alert-level correspondence, and we report it as agreement rather than as
confirmation. The record also sees the alerts at only 3 of 10 window--order
cells: at 0.55 (canonical), 0.62 (canonical), 0.62 (first-flow) the issued alerts sit in the \emph{feasible prefix} --- the first
minutes of a deployment window --- which predates the exercise's first submission, so at most
$9\%$ of them overlap it in time at all; at 0.85 (canonical), 0.85 (first-flow) they do overlap but only $2$ of $17$ tasks with a
timed step declare a segment or file a compromise report. Four further limits. \emph{(i)} At the
$86{,}400$\,s bucket both nulls tie the observation exactly: every submission shares one daily bucket,
so no permutation and no rotation can move a step between alerts (the spatial filter is \emph{not}
vacuous there --- at $0.70$/first-flow it admits $12$ of $30$ temporally overlapping alerts --- it is
the nulls that lose their power); only sub-daily widths inform. \emph{(ii)} The whole positive result
rests on six coarse segment labels: across every row the compromise-IP conjunct adds \emph{no} alert
the segment conjunct did not already admit. \emph{(iii)} Every attributed alert sits on a
labelled-malicious episode and the labels came from this same exercise, so nothing here is
label-independent evidence that an alert is a true positive. \emph{(iv)} A submission is when the red
team \emph{reported}: the lag, estimable from the compromise reports' own timestamps, has median
170\,s (IQR 115--279\,s), small against every bucket width.
